% Preámbulo modular para presentaciones Beamer (Modelación Estadística)
% Diseño y paleta institucional del Tecnológico de Monterrey

\documentclass[aspectratio=169,xcolor={dvipsnames,table}]{beamer}

\usepackage[utf8]{inputenc}
\usepackage[spanish,mexico]{babel}
\usepackage{amsmath,amssymb,amsfonts,mathtools}
\usepackage{graphicx}
\usepackage{listings}
\usepackage{booktabs}
\usepackage{tikz}
\usetikzlibrary{matrix,arrows,decorations.pathmorphing,shapes.geometric,positioning}

% Paleta Institucional Tecnológico de Monterrey
\definecolor{TecAzul}{HTML}{0039A6}       % Azul TEC: color primario institucional
\definecolor{TecAzulOscuro}{HTML}{1A2E51} % Azul marino oscuro
\definecolor{TecRojo}{HTML}{EC2661}       % Rojo de acento (paleta miTec)
\definecolor{TecGris}{HTML}{646464}       % Gris neutro para comentarios y subtítulos
\definecolor{TecFondoBloque}{HTML}{F4F6F9}% Fondo suave para bloques

% Configuración de Tema Beamer
\usetheme{Boadilla}
\usecolortheme[named=TecAzulOscuro]{structure}
\setbeamercolor{alerted text}{fg=TecRojo}
\setbeamercolor{palette primary}{bg=TecAzulOscuro,fg=white}
\setbeamercolor{palette secondary}{bg=TecAzul,fg=white}
\setbeamercolor{palette tertiary}{bg=TecAzulOscuro!80!black,fg=white}
\setbeamercolor{title}{fg=TecAzulOscuro,bg=TecFondoBloque}
\setbeamercolor{frametitle}{fg=TecAzulOscuro,bg=TecFondoBloque}

% Bloques personalizados elegantes
\setbeamercolor{block title}{bg=TecAzulOscuro,fg=white}
\setbeamercolor{block body}{bg=TecFondoBloque,fg=black}
\setbeamercolor{block title alerted}{bg=TecRojo,fg=white}
\setbeamercolor{block body alerted}{bg=TecRojo!8,fg=black}
\setbeamercolor{block title example}{bg=TecAzul,fg=white}
\setbeamercolor{block body example}{bg=TecAzul!8,fg=black}

% Redondear bordes y añadir sombra sutil
\setbeamertemplate{blocks}[rounded][shadow=true]
\setbeamertemplate{navigation symbols}{}

% Pie de página limpio con número de diapositiva
\setbeamertemplate{footline}{
  \leavevmode%
  \hbox{%
  \begin{beamercolorbox}[wd=.7\paperwidth,ht=2.25ex,dp=1ex,left]{author in head/foot}%
    \hspace*{2ex}\insertshortauthor~~(\insertshortinstitute) \hfill \insertshorttitle
  \end{beamercolorbox}%
  \begin{beamercolorbox}[wd=.3\paperwidth,ht=2.25ex,dp=1ex,right]{date in head/foot}%
    \insertshortdate{}\hspace*{2em}
    \insertframenumber{} / \inserttotalframenumber\hspace*{2ex}
  \end{beamercolorbox}}%
  \vskip0pt%
}

% Configuración de Listings de Python para visualización pedagógica de código
\lstset{
    language=Python,
    basicstyle=\ttfamily\footnotesize,
    keywordstyle=\color{TecAzulOscuro}\bfseries,
    stringstyle=\color{TecRojo},
    commentstyle=\color{TecGris}\itshape,
    showstringspaces=false,
    breaklines=true,
    frame=lines,
    rulecolor=\color{TecAzulOscuro!30},
    backgroundcolor=\color{TecFondoBloque},
    aboveskip=0.5em,
    belowskip=0.5em,
    literate={á}{{\'a}}1 {é}{{\'e}}1 {í}{{\'i}}1 {ó}{{\'o}}1 {ú}{{\'u}}1
             {Á}{{\'A}}1 {É}{{\'E}}1 {Í}{{\'I}}1 {Ó}{{\'O}}1 {Ú}{{\'U}}1
             {ñ}{{\~n}}1 {Ñ}{{\~N}}1 {¿}{{?`}}1 {¡}{{!`}}1
}

% Metadatos institucionales por defecto
\title[Modelación Estadística]{Modelación Estadística para Ciencia de Datos e Ingeniería}
\author[J. Castillo Colmenares]{Juliho David Castillo Colmenares}
\institute[Tec de Monterrey]{Tecnológico de Monterrey \\ Escuela de Ingeniería y Ciencias}
\date{\today}

% Comandos y macros estadísticas compartidas para las presentaciones Beamer (Metropolis)

% Conjuntos numéricos básicos
\newcommand{\R}{\mathbb{R}}
\newcommand{\N}{\mathbb{N}}
\newcommand{\Q}{\mathbb{Q}}
\newcommand{\Z}{\mathbb{Z}}
\newcommand{\set}[1]{\left\{ #1 \right\}}
\newcommand{\abs}[1]{\left|#1\right|}
\newcommand{\norm}[1]{\left\| #1 \right\|}

% Comandos estadísticos y probabilísticos del curso
\newcommand{\Var}{\operatorname{Var}}
\newcommand{\cov}{\operatorname{Cov}}
\newcommand{\corr}{\rho}
\newcommand{\comb}[2]{\begin{pmatrix} #1 \\ #2 \end{pmatrix}}
\newcommand{\s}{\sigma}
\newcommand{\particion}{\mathfrak{P}}
\newcommand{\card}[1]{\#\left( #1 \right)}
\newcommand{\seq}[3]{\set{#1}_{#2}^{#3}}
\newcommand{\E}{\mathbb{E}}
\newcommand{\Prob}{\mathbb{P}}

% Entornos pedagógicos y cajas en español académico natural
\newenvironment{discusion}{%
  \begin{alertblock}{Pregunta de Discusión}%
}{%
  \end{alertblock}%
}

\newenvironment{reflexion}{%
  \begin{alertblock}{Reflexión para el Aula}%
}{%
  \end{alertblock}%
}

\newenvironment{paradoja}{%
  \begin{alertblock}{Paradoja de Exploración}%
}{%
  \end{alertblock}%
}

\newenvironment{motivacion}{%
  \begin{exampleblock}{Motivación: Ciencia de Datos en la Práctica}%
}{%
  \end{exampleblock}%
}

\newenvironment{retoclase}{%
  \begin{block}{Reto Rápido en Equipo}%
}{%
  \end{block}%
}


\title[Probability Introduction]{Section 2.1: Introduction to Probability}
\subtitle{Mathematical Modeling of Randomness and Uncertainty}
\author[J. Castillo Colmenares]{Juliho Castillo Colmenares}
\institute{Tecnológico de Monterrey}
\date{}

\begin{document}

% Slide 1: Title Page
\begin{frame}[plain]
  \titlepage
\end{frame}

% Slide 2: Roadmap and Position Indicator
\begin{frame}{Roadmap — Chapter 02: Probability Theory}
  Where are we in our journey through probability theory? \pause
  \begin{itemize}
    \item \color{TecRojo}\textbf{02.01 Introduction to probability \emph{(Today)}}\color{black} \pause
    \item \textbf{02.02 Set notation and partitions} \emph{(Next session)} \pause
    \item \textbf{02.03 Foundations and probability axioms} \pause
    \item \textbf{02.04 Conditional probability and independence} \pause
    \item \textbf{02.05 Bayes' Theorem} \pause
    \item \textbf{02.06 Random sampling}
  \end{itemize}
  \vspace{0.8em}
  \begin{block}{Section 2.1 Objective}
    Understand why we must model random phenomena where classical determinism fails, and establish the deductive/inductive bridge between Probability and Data Science.
  \end{block}
\end{frame}

% Slide 3: Motivation / Trigger Question
\begin{frame}{Why Can't We Predict Everything Exactly?}
  \begin{motivacion}
    If we know gravity, the mass of a rocket, and engine thrust, classical physics lets us predict its trajectory with millimeter accuracy. However, can we predict the exact minute an \textbf{Amazon} user will click "Buy"? Or whether a patient will respond positively to a new cancer therapy?
  \end{motivacion}
  
  \vspace{0.6em}
  \pause
  \begin{itemize}
    \item In complex biological, social, and technological systems, the number of unobservable or uncontrollable variables is immense. \pause
    \item \textbf{Probability Theory} provides the exact mathematical language required to reason, quantify, and make decisions under this unavoidable uncertainty.
  \end{itemize}
\end{frame}

% Slide 4: Deterministic vs. Random Phenomena
\begin{frame}{Deterministic vs. Random Phenomena}
  \begin{columns}[T]
    \begin{column}{0.48\textwidth}
      {\large \color{TecAzulOscuro}\textbf{Deterministic Phenomenon}}
      \vspace{0.3em}
      \begin{itemize}
        \item \textbf{Definition:} Under identical initial conditions, the outcome is always precisely the same and predictable with certainty. \pause
        \item \textbf{Examples:} Free fall in a vacuum ($d = \frac{1}{2}gt^2$), area of a circle ($\pi r^2$), fixed-rate compound interest. \pause
        \item \textbf{Uncertainty:} $0\%$.
      \end{itemize}
    \end{column}
    
    \begin{column}{0.48\textwidth}
      {\large \color{TecRojo}\textbf{Random Phenomenon}}
      \vspace{0.3em}
      \begin{itemize}
        \item \textbf{Definition:} Its exact outcome cannot be predicted beforehand with certainty, even when replicating initial conditions. \pause
        \item \textbf{Examples:} Web server latency, stock market fluctuations, rolling a die. \pause
        \item \textbf{Uncertainty:} Governed by stable statistical regularities over the long run.
      \end{itemize}
    \end{column}
  \end{columns}
\end{frame}

% Slide 5: The Probability vs. Statistics Bridge
\begin{frame}{The Dual Bridge: Probability vs. Statistics}
  Probability and statistics are two sides of the same scientific reasoning coin:
  
  \vspace{0.2em}
  \begin{block}{The Deductive Engine vs. The Inductive Telescope}
    \begin{itemize}
      \item \textbf{Probability (Deduction):} Starts from a known \emph{theoretical population model} and predicts the likelihood of the sample data we might observe. \pause
      \item \textbf{Inferential Statistics (Induction):} Starts from observed \emph{sample data} and uses the laws of probability to infer or estimate the underlying population model.
    \end{itemize}
  \end{block}
  
  \vspace{0.2em}
  \pause
  \begin{center}
    \small \textbf{Without probability theory, inferential statistics would lack the rigorous mathematical foundation needed to quantify margins of error.}
  \end{center}
\end{frame}

% Slide 6: Applications in Science & Engineering
\begin{frame}{Applications in the Data Age}
  Probability is the computational engine driving modern technologies:
  \vspace{0.4em}
  \begin{itemize}
    \item \textbf{Artificial Intelligence \& Machine Learning:} Probabilistic neural networks, Naive Bayes classifiers, and reinforcement learning. \pause
    \item \textbf{Site Reliability Engineering (SRE):} Modeling server time-to-failure, fault tolerance, and queuing theory in cloud infrastructure. \pause
    \item \textbf{Quantitative Finance:} Black-Scholes option pricing, Value at Risk ($VaR$), and stochastic portfolio optimization. \pause
    \item \textbf{Bioinformatics \& Genetics:} Genome decoding, evolutionary mutation rates, and double-blind clinical trials.
  \end{itemize}
\end{frame}

% Slide 7: Historical Evolution
\begin{frame}{Historical Evolution of Probability}
  \begin{itemize}
    \item \textbf{17th Century (The Birth):} In $1654$, an epistolary correspondence between \textbf{Blaise Pascal} and \textbf{Pierre de Fermat} regarding fair division of stakes in interrupted gambling games established the algebraic foundation of mathematical probability. \pause
    \item \textbf{18th \& 19th Centuries (Consolidation):} \textbf{Jacob Bernoulli} proved the first Law of Large Numbers ($1713$), and \textbf{Carl Friedrich Gauss} characterized the normal distribution of measurement errors ($1809$). \pause
    \item \textbf{20th Century (Modern Axiomatization):} In $1933$, \textbf{Andrey Kolmogorov} published \emph{Foundations of the Theory of Probability}, unifying probability with Lebesgue measure theory using just three indisputable axioms.
  \end{itemize}
\end{frame}

% Slide 8: Chapter 02 Structural Overview
\begin{frame}{Chapter Structure: What We Will Build}
  In the upcoming sections, we will master the formal tools of random modeling:
  \vspace{0.6em}
  \begin{enumerate}
    \item \textbf{Algebra of Events (Section 2.2):} Using sets, Venn diagrams, and partitions to divide and count the universe of possibilities without double counting. \pause
    \item \textbf{Kolmogorov's Axioms (Section 2.3):} The three fundamental rules from which every probabilistic formula is derived. \pause
    \item \textbf{Conditional Probability (Section 2.4):} Updating and adjusting probabilities when partial information or prior observations become available. \pause
    \item \textbf{Bayes' Theorem (Section 2.5):} The core engine of machine learning that updates prior beliefs in light of new evidence.
  \end{enumerate}
\end{frame}

% Slide 9: Summary & Connection
\begin{frame}{Summary and Connection to Section 2.2}
  \begin{block}{Key Takeaways}
    \begin{itemize}
      \item Random phenomena possess inherent variability whose long-term structure is governed by precise probabilistic laws. \pause
      \item Probability is the mathematical bedrock of both inferential statistics and modern Machine Learning algorithms.
    \end{itemize}
  \end{block}
  
  \vspace{0.8em}
  \pause
  {\large \color{TecAzulOscuro}\textbf{Next Section: 02.02 Set Notation and Partitions}}
  \vspace{0.3em}
  \begin{itemize}
    \item To compute probabilities without ambiguity, we must first learn to define the universe of outcomes algebraically using \textbf{Sets} and decompose it like a puzzle using \textbf{Partitions}.
  \end{itemize}
\end{frame}

\end{document}
